\begin{scanpage}{0187}{169}
curvature, a sharp bend, a point, a bulge, a hole, etc.). One should note that, in the case of electromagnetic wave scattering by a system of separate bodies, the separable part of the field is due not only to the surface's distortion, but also to the diffraction interaction of the bodies.

It is necessary, however, to keep in mind that it is possible to realize the indicated fringing field distribution not in an arbitrary observation direction, but only in a direction for which the condition $\Sigma^0=-\overline\Sigma^0$ is fulfilled---for example, in the direction towards the source.

Consideration of the nonuniform part of the current also enables one to explain the reflected wave depolarization which we will investigate in the following section.

Figure 71 presents the results of measurements$^{(2)}$ and calculations of the effective scattering surface
\[
\sigma^+=\pi a^2|\Sigma_+|^2=\frac14\pi a^2|\Sigma^1+\overline\Sigma^1|^2,\tag{26.19}
\]
which is dependent upon the nonuniform part of the current excited by a plane electromagnetic wave on a disk. The disk's diameter equals $2a=5\lambda/\pi$ ($\lambda$ is the wavelength). The calculations were performed with consideration of the secondary diffraction on the basis of the approximation equations for the functions $\Sigma$ and $\overline\Sigma$ which were derived in \S~24. Since it is difficult to prepare a thin disk with a sufficiently flat surface, the measurements were performed with an obtuse cone close to the shape of a disk and having a height approximately equal to one tenth of the diameter.

As is seen from Figure 71, the theoretical and experimental curves are fairly close together. A certain divergence between them, especially in the region of $\gamma$ values close to $90^\circ$, may evidently be explained both by the model's conical shape and also by the
\par\bigskip\hrule\smallskip
$(2)$Footnote appears on page 174.
\end{scanpage}
